\documentclass[10pt]{article}
\usepackage[hmargin=1in,vmargin=0.9in]{geometry}
\usepackage{amsmath,amssymb,amsthm}
\usepackage{booktabs}
\usepackage{microtype}
\usepackage[hidelinks]{hyperref}
\hypersetup{pdftitle={An effective finiteness theorem for an exceptional set of Erdos and Selfridge},
  pdfauthor={Haoyu Chen}, pdfsubject={Erdos Problem 889}}

\numberwithin{equation}{section}
\newtheorem{theorem}{Theorem}[section]
\newtheorem{lemma}[theorem]{Lemma}
\newtheorem{proposition}[theorem]{Proposition}
\newtheorem{corollary}[theorem]{Corollary}
\theoremstyle{definition}
\newtheorem{definition}[theorem]{Definition}
\newtheorem{remark}[theorem]{Remark}

% named statements (Theorem A, Corollary A1, ...)
\newcommand{\stmtname}{}
\theoremstyle{plain}
\newtheorem*{named}{\stmtname}
\newenvironment{namedthm}[1]{\renewcommand{\stmtname}{#1}\begin{named}}{\end{named}}

\makeatletter
\renewcommand\section{\@startsection{section}{1}{\z@}{-2.6ex \@plus -.8ex \@minus -.2ex}{1.3ex \@plus .2ex}{\normalfont\large\bfseries}}
\renewcommand\subsection{\@startsection{subsection}{2}{\z@}{-2.2ex \@plus -.6ex \@minus -.2ex}{1ex \@plus .2ex}{\normalfont\normalsize\bfseries}}
\makeatother

\newcommand{\cP}{\mathcal{P}}
\newcommand{\cY}{\mathcal{Y}}
\newcommand{\Bbar}{\overline{B}}
\newcommand{\unif}{\mathrm{u}}
\newcommand{\eps}{\varepsilon}

\title{An effective finiteness theorem for an exceptional set\\ of Erd\H{o}s and Selfridge}
\author{Haoyu Chen\\ Independent researcher}
\date{September 25, 2026}

\begin{document}
\maketitle

\begin{abstract}
For integers $n\ge 0$ and $k\ge 1$ let $v(n,k)$ be the number of prime factors of $n+k$ that divide none
of $n,n+1,\dots,n+k-1$, equivalently the number of primes $p>k$ dividing $n+k$, and let
$v_1(n)=\sup_{k\ge 1}v(n,k)$. Erd\H{o}s and Selfridge (1967) expected $v_1(n)=1$ to have only finitely many
solutions but could not prove it. We prove that every $n\ge N_0$, where
$\log N_0=e^{45.28}\approx 4.62\cdot 10^{19}$, has some $k\le 10\log n$ with $v(n,k)\ge 2$. Hence every
solution of $v_1(n)=1$ is smaller than $N_0$. The proof combines an elementary counting lemma with
Matveev's explicit lower bound for linear forms in logarithms; every numerical inequality is certified by
interval arithmetic. The same argument gives, for all integers $l\ge1$ and $n$ with $\log\log n\ge 48.82$ and
$\log n\ge l$, some $k\in[l,10\,l\log n]$ with $v(n,k)\ge 2$. The qualitative finiteness also follows from
a 1981 theorem of Langevin combined with an elementary lemma.
\end{abstract}

\noindent\textit{MSC 2020:} 11N25, 11J86.\quad
\textit{Keywords:} prime factors of consecutive integers; smooth numbers; linear forms in logarithms;
Erd\H{o}s problems.

\section{Introduction}

Erd\H{o}s and Selfridge \cite{ErSe67} considered, for integers $n\ge 0$ and $k\ge 0$, the number $v(n;k)$
of prime factors of $n+k$ that divide none of $n,n+1,\dots,n+k-1$, and the quantities
\[
v_0(n)=\max_{k\ge 0}v(n;k),\qquad v_l(n)=\max_{k\ge l}v(n;k).
\]
For positive integers $n$, they showed that $v_0(n)>1$ except when $n=1,2,3,4,7,8,16$. They wrote: ``It seems certain that
$\lim v_l(n)=\infty$ for every $l$, but unfortunately we have not even been able to prove that $v_1(n)=1$ has
only a finite number of solutions, though this certainly must be true. Probably the greatest $n$ for which
$v_1(n)=1$ is $n=330$.'' They listed the $29$ values $n<2500$ with $v_1(n)=1$:
$n=1$--$4$, $6$--$8$, $10$, $12$, $15$, $16$, $18$, $22$, $24$, $26$, $30$, $36$, $42$, $46$, $48$, $60$, $70$,
$78$, $80$, $96$, $120$, $190$, $222$, $330$.
Guy's \emph{Unsolved Problems in Number Theory} \cite[B27]{Gu04} records the question in its third edition
(2004), repeating that Erd\H{o}s and Selfridge were unable to prove that $v_1(n)=1$ has only finitely many
solutions. The problem is Problem \#889 of the Erd\H{o}s problems database \cite{EP889}, and it is formalised
in the Formal Conjectures repository \cite{FC}, where the variant \texttt{erdos\_889.variants.v1\_eq\_1\_finite}
asks whether $\{n: v_1(n)=1\}$ is finite.

Erd\H{o}s and Selfridge also introduced a variant $V(n;k)$, the number of primes $p\mid n+k$ whose exact
prime-power component $p^{\alpha}\,\|\,n+k$ divides none of $n,\dots,n+k-1$, and asked the same question for
$V_1(n)=1$ (perhaps with largest solution $80$). Kitamura \cite{Ki26} has given a Lean~4 proof that
$\{n:V_1(n)=1\}$ is finite (repository dated 22 September 2026); its documentation states that it does not
address the lowercase variant. Since $p\nmid n+i$ implies $p^\alpha\nmid n+i$, one has $v(n;k)\le V(n;k)$, so
$V_1(n)\ge v_1(n)$ and finiteness of $\{n:v_1(n)=1\}$ implies finiteness of $\{n:V_1(n)=1\}$, not conversely.
(In particular, Theorem~A below also shows that every solution of $V_1(n)=1$ is smaller than $N_0$.)

\paragraph{What this note contains.} An explicit threshold $N_0$ with $\log\log N_0=45.28$ beyond which
$v(n,k)\ge2$ for some $k\le10\log n$ (Theorem~A), and an explicit version uniform in $l$. An elementary lemma
(Lemma~\ref{lem:R}) shows that if $v(n,k)\le1$ for all $k\le y$, then two of $n+1,\dots,n+y$ have the form
$Wq^{e}$ with distinct primes $q$ and small cofactors $W$; Matveev's bound \cite{Ma00} applied to their quotient
gives $\log n< K(\log\log n)^5$ with an explicit $K$.

\paragraph{What is not new.} The qualitative statement that $\{n:v_1(n)=1\}$ is finite, with an effective
(but not computed) bound, is a corollary of a published theorem: Langevin's 1981 bound for the number of
$k$-smooth integers among $n+1,\dots,n+k$ \cite[Corollaire~1, (8)]{La81}, combined with a ten-line elementary
lemma. We write out this deduction in Section~\ref{sec:langevin}; the same route also gives the qualitative
uniform-in-$l$ statement. The method is also classical: Baker-type lower bounds applied to smooth numbers in
blocks of consecutive integers go back to Tijdeman \cite{Ti73}, Ramachandra, Shorey and Tijdeman
\cite{RST75,RST76} and Langevin \cite{La81}, and the device of assigning each prime to the term in which it occurs
to the highest power is of Sylvester--Erd\H{o}s type. What this note adds is quantitative: an explicit threshold
$\log\log N_0=45.28$ and an explicit uniform-in-$l$ threshold $48.82$. The threshold is far beyond any
computation, so the conjecture that $330$ is the largest solution of $v_1(n)=1$ remains open, as does the main
problem $v_0(n)\to\infty$.

\section{Notation and statements}\label{sec:statements}

Throughout, $\log$ is the natural logarithm and $p,q$ denote primes. For an integer $m\ge2$, $P(m)$ is its
largest prime factor, and $P(1)=1$; $\omega(m)$ is the number of distinct prime factors of $m$; $v_p$ is the
$p$-adic valuation; and $\pi(x)=\#\{p\le x\}$ for real $x$. Following \cite{ErSe67} and \cite{FC}, for integers
$n\ge0$, $k\ge1$ and $l\ge1$ put
\[
v(n,k)=\#\{p : p\mid n+k,\ p\nmid n+i\ \text{for all } 0\le i\le k-1\},\qquad
v_l(n)=\sup_{k\ge l}v(n,k)\in\mathbb{N}\cup\{\infty\}.
\]

\begin{lemma}\label{lem:v}
For integers $n\ge0$ and $k\ge1$, $v(n,k)=\#\{p : p\mid n+k,\ p>k\}$.
\end{lemma}

\begin{proof}
Let $p\mid n+k$. If $p\le k$, then $i:=k-p$ satisfies $0\le i\le k-1$ and $p\mid (n+k)-p=n+i$, so $p$ is not
counted. If $p>k$ and $0\le i\le k-1$, then $(n+k)-(n+i)=k-i\in[1,k]$ is not divisible by $p$, so $p\nmid n+i$
and $p$ is counted. Thus $v(n,k)\le1$ means that $n+k$ has at most one prime factor greater than $k$.
\end{proof}

\begin{namedthm}{Theorem A}
Put $\ell_1:=45.28$ and $N_0:=\exp(\exp(\ell_1))$, so that $\log N_0=e^{45.28}\approx4.6223\cdot10^{19}$ $($to five
figures$)$. For every integer $n\ge N_0$ there is an integer $k$ with $1\le k\le 10\log n$ and $v(n,k)\ge 2$.
\end{namedthm}

\begin{namedthm}{Corollary A1}
Every $n$ with $v_1(n)=1$ satisfies $n<N_0$; in particular $\{n: v_1(n)=1\}$ is finite.
\end{namedthm}

\begin{proof}
If $n\ge N_0$, Theorem~A gives $k\ge1$ with $v(n,k)\ge2$, so $v_1(n)\ge2$ and $v_1(n)\ne1$.
\end{proof}

In the terms of \cite{FC}, Corollary~A1 is a paper proof, not a Lean proof, that
\texttt{Erdos889.\allowbreak erdos\_889.\allowbreak variants.\allowbreak v1\_eq\_1\_finite} holds with answer
\texttt{True}. It depends on the
published results of Matveev, Rosser--Schoenfeld and Robin (Sections~\ref{sec:baker}--\ref{sec:numerics}). The
qualitative finiteness also follows from Langevin \cite{La81} with Lemma~\ref{lem:R}(i) (Section~\ref{sec:langevin}).

\begin{namedthm}{Theorem A$_l$ (fixed $l$)}
Let $1\le l\le 20$ and let $\ell_1(l)$ be as in Table~\ref{tab:perl} $($so $\ell_1(1)=45.28$, $\ell_1(2)=45.35$,
\dots, $\ell_1(20)=45.58)$. Every integer $n$ with $\log\log n\ge\ell_1(l)$ has an integer $k$ with
$l\le k\le 10\,l\log n$ and $v(n,k)\ge 2$.
\end{namedthm}

\begin{namedthm}{Theorem A$_{\mathrm{unif}}$ (uniform in $l$)}
Let $l\ge1$ and $n$ be integers with $\log\log n\ge 48.82$ and $\log n\ge l$. Then there is an integer $k$ with
$l\le k\le 10\,l\log n$ and $v(n,k)\ge2$. Equivalently, this holds for every $n\ge N(l)$, where
\[
N(l):=\max\bigl(\exp(e^{48.82}),\ \lceil e^{l}\rceil\bigr),\qquad e^{48.82}\approx1.593\cdot10^{21}.
\]
\end{namedthm}

\begin{namedthm}{Corollary A$_l$}
For every $l\ge1$ and every $n\ge N(l)$, $v_l(n)\ge2$.
\end{namedthm}

Corollary~A$_l$ is the level-$2$ case of \texttt{erdos\_889.variants.general} in \cite{FC}; that variant asks
for $v_l(n)\to\infty$, which is not proved here. Section~\ref{sec:open} lists other questions that this note
does not address.

\section{The elementary reduction}\label{sec:reduction}

\paragraph{Setting.} Fix integers $l\ge1$ and $n\ge2$ and a real number $y$ with $Y:=\lfloor y\rfloor\ge l+2$.
Put $\Pi_l(n):=n(n+1)\cdots(n+l-1)$ and
\[
\omega_{l,y}(n):=\#\{p\le y: p\mid \Pi_l(n)\},\qquad \cP:=\{p\le y: p\nmid\Pi_l(n)\},\qquad
|\cP|=\pi(y)-\omega_{l,y}(n).
\]
The hypothesis is
\begin{equation}\tag{H}
v(n,k)\le1\quad\text{for every integer } k \text{ with } l\le k\le y.
\end{equation}
Define
\[
A:=\pi(y)-2\,\omega_{l,y}(n)-\frac{2\log\bigl((Y-l-1)!\bigr)}{\log n},\qquad\text{and, when } A>1,\qquad
T:=\frac{\log\bigl((Y-l-1)!\bigr)}{A-1}.
\]

\begin{lemma}[Lemma R]\label{lem:R}
Assume \textup{(H)}.
\begin{enumerate}
\item[(i)] For $p\in\cP$ let $k_0(p):=\min\{k\ge l: p\mid n+k\}$. Then $l\le k_0(p)\le p-1\le Y-1$ and
$P(n+k_0(p))=p$. The map $p\mapsto k_0(p)$ is injective. Hence $J:=k_0(\cP)\subseteq\{l,l+1,\dots,Y-1\}$ has
$|J|=\pi(y)-\omega_{l,y}(n)$, and for $j\in J$ every prime factor of $n+j$ is at most $Y$.
\item[(ii)] Assume $J\ne\emptyset$. For every prime $p\le Y$ fix $j_p\in J$ with
$v_p(n+j_p)=\max_{j\in J}v_p(n+j)$. For $j\in J$ let
\[
U_j:=\prod_{p\le Y,\ j_p=j}p^{v_p(n+j)},\qquad W_j:=\prod_{p\le Y,\ j_p\ne j}p^{v_p(n+j)}.
\]
Then $n+j=U_jW_j$ for each $j\in J$, and $\prod_{j\in J}W_j\le (Y-l-1)!$.
\item[(iii)] Let $s(j):=\omega(U_j)$, $J_0:=\{j\in J:s(j)=0\}$ and $J_1:=\{j\in J:s(j)=1\}$. Then
$|J_0|\le\log((Y-l-1)!)/\log n$ and $|J_1|\ge\pi(y)-2\,\omega_{l,y}(n)-2|J_0|\ge A$.
\item[(iv)] For $j\in J_1$, $U_j=q_j^{e_j}$ with $q_j$ a prime $\le Y$ and $e_j\ge1$. So $n+j=W_jq_j^{e_j}$ with
$q_j\nmid W_j$. The map $j\mapsto q_j$ is injective on $J_1$.
\item[(v)] If $A>1$, then $J\ne\emptyset$, and there are $j_1\ne j_2$ in $J_1$ with $\log W_{j_1}\le T$ and
$\log W_{j_2}\le T$.
\end{enumerate}
\end{lemma}

\begin{proof}
(i) Let $p\in\cP$. Among the $p$ consecutive integers $n,\dots,n+p-1$ there is a multiple $n+i$ of $p$, with
$0\le i\le p-1$. If $p\le l$, this $n+i$ divides $\Pi_l(n)$, contradicting $p\nmid\Pi_l(n)$; so $p\ge l+1$.
Since $p\nmid n+i$ for $i<l$, the multiple has $l\le i\le p-1$; hence $l\le k_0(p)\le p-1$. As $p\le y$ is an
integer, $p\le Y$, so $k_0(p)\le Y-1<y$ and (H) gives $v(n,k_0(p))\le1$. By Lemma~\ref{lem:v} (note
$k_0(p)\ge1$), $n+k_0(p)$ has at most one prime factor greater than $k_0(p)$. The prime $p$ is such a factor, so
every other prime factor of $n+k_0(p)$ is at most $k_0(p)<p$, and $P(n+k_0(p))=p\le Y$. If $k_0(p)=k_0(p')$, then
$p=P(n+k_0(p))=p'$; so the map is injective and $|J|=|\cP|$.

(ii) By (i), $n+j=\prod_{p\le Y}p^{v_p(n+j)}$ for $j\in J$; splitting the primes $p\le Y$ according to whether
$j_p=j$ gives $n+j=U_jW_j$. Fix a prime $p\le Y$ and $j\in J$ with $j\ne j_p$. Since
$v_p(n+j)\le v_p(n+j_p)$, the power $p^{v_p(n+j)}$ divides both $n+j$ and $n+j_p$, hence divides
$j-j_p\ne0$, and $v_p(n+j)\le v_p(|j-j_p|)$. Summing over $j\in J\setminus\{j_p\}$, a subset of
$\{l,\dots,Y-1\}\setminus\{j_p\}$, and using that all terms are nonnegative,
\[
\sum_{j\in J,\ j\ne j_p}v_p(n+j)\le\sum_{\substack{i=l\\ i\ne j_p}}^{Y-1}v_p(|i-j_p|)
=v_p\bigl((j_p-l)!\,(Y-1-j_p)!\bigr)\le v_p\bigl((Y-l-1)!\bigr),
\]
because $a!\,b!$ divides $(a+b)!$ and here $a+b=Y-l-1$. Therefore
\[
\prod_{j\in J}W_j=\prod_{p\le Y}p^{\sum_{j\in J,\,j\ne j_p}v_p(n+j)}\le\prod_{p\le Y}p^{v_p((Y-l-1)!)}=(Y-l-1)!,
\]
the last equality because every prime factor of $(Y-l-1)!$ is at most $Y$.

(iii) For $j\in J_0$ we have $U_j=1$, so $W_j=n+j>n$. All $W_j\ge1$, so by (ii)
$|J_0|\log n\le\sum_{j\in J_0}\log W_j\le\sum_{j\in J}\log W_j\le\log((Y-l-1)!)$. Next,
$s(j)=\#\{p\le Y: j_p=j,\ v_p(n+j)\ge1\}$. Every prime $p\le Y$ has exactly one $j_p$, so these sets are
pairwise disjoint subsets of the primes up to $Y$, and $\sum_{j\in J}s(j)\le\pi(Y)=\pi(y)$. With
$J_{\ge2}:=J\setminus(J_0\cup J_1)$ this gives $|J_1|+2|J_{\ge2}|\le\pi(y)$, while
$|J_0|+|J_1|+|J_{\ge2}|=|J|=\pi(y)-\omega_{l,y}(n)$. Eliminating $|J_{\ge2}|$ gives
$|J_1|\ge\pi(y)-2\omega_{l,y}(n)-2|J_0|$, which is at least $A$ by the bound for $|J_0|$.

(iv) If $s(j)=1$, there is exactly one prime $q=q_j\le Y$ with $j_q=j$ and $v_q(n+j)\ge1$; the other primes $p$
with $j_p=j$ contribute $p^0=1$. So $U_j=q_j^{e_j}$ with $e_j:=v_{q_j}(n+j)\ge1$, and $q_j\nmid W_j$ because
$W_j$ contains only primes $p$ with $j_p\ne j$. Finally $q_j$ determines $j=j_{q_j}$.

(v) If $J=\emptyset$, then $\omega_{l,y}(n)=\pi(y)$ and $A\le-\pi(y)<1$; so $A>1$ forces $J\ne\emptyset$. By (iii),
$r:=|J_1|\ge A>1$, so the integer $r$ is at least $2$. List the numbers $\log W_j$ $(j\in J_1)$ as
$w_1\le w_2\le\dots\le w_r$ and let $j_1,j_2$ be indices giving $w_1,w_2$. All $w_i\ge0$, so by (ii)
$(r-1)w_2\le w_2+\dots+w_r\le\sum_{j\in J}\log W_j\le\log((Y-l-1)!)$. Since $r-1\ge A-1>0$, we get
$w_1\le w_2\le T$.
\end{proof}

\section{The Baker step}\label{sec:baker}

\subsection{Matveev's theorem}
We use Corollary~2.3 of Matveev \cite{Ma00} (Russian original p.~127, English translation p.~1219), in his
notation and equation numbering, except that his number $n$ of logarithms is called $t$ here. Let $\mathbb{K}\subset\mathbb{C}$ be an
algebraic number field of degree $D$ over $\mathbb{Q}$; put $\varkappa=1$ if $\mathbb{K}\subseteq\mathbb{R}$ and
$\varkappa=2$ otherwise. Let $\alpha_1,\dots,\alpha_t\in\mathbb{K}^*$ with absolute logarithmic heights
$h(\alpha_j)$, let $\ln\alpha_1,\dots,\ln\alpha_t$ be arbitrary fixed \emph{non-zero} values of the logarithms,
let $b_1,\dots,b_t\in\mathbb{Z}$, and put
\begin{gather*}
\Lambda=b_1\ln\alpha_1+\dots+b_t\ln\alpha_t\quad(1.1),\qquad
B=\max\Bigl\{1,\max_{1\le j\le t}\frac{|b_j|A_j}{A_t}\Bigr\}\quad(1.3),\qquad
B^*=\max_{1\le j\le t}|b_j|\quad(1.4),\\
A_j\ge\max\{D\,h(\alpha_j),\ |\ln\alpha_j|,\ 0.16\}\quad(2.4),\qquad \Omega=A_1\cdots A_t.
\end{gather*}

\begin{theorem}[{Matveev \cite[Corollary~2.3]{Ma00}}]\label{thm:matveev}
If $\Lambda\ne0$, the $A_j$ satisfy $(2.4)$ and $B$ is given by $(1.3)$, then
\[
\ln|\Lambda|>-C_1(t)\,D^2\,\Omega\,\ln(eD)\,\ln(eB),\qquad
C_1(t)=C_1(t,\varkappa)=\min\Bigl\{\frac1\varkappa\Bigl(\frac12et\Bigr)^{\varkappa}30^{t+3}t^{3.5},\ 2^{6t+20}\Bigr\},
\]
where $B$ may be replaced by $B^*$.
\end{theorem}

We use it only with $\mathbb{K}=\mathbb{Q}\subset\mathbb{R}$, so $D=1$, $\varkappa=1$ and $\ln(eD)=1$, with the
real logarithms. For a rational number $a/b$ in lowest terms, $h(a/b)=\log\max(|a|,|b|)$. The constants are
\begin{align*}
C_1(3)&=\min\{\tfrac{3e}{2}\,30^6\,3^{3.5},\,2^{38}\}=\tfrac{3e}{2}\,30^6\,3^{3.5}=1.39007\ldots\cdot10^{11},\\
C_1(2)&=\min\{e\,30^5\,2^{3.5},\,2^{32}\}=e\,30^5\,2^{3.5}=7.47318\ldots\cdot10^{8}.
\end{align*}
The frequently quoted simplified form $1.4\cdot30^{t+3}t^{4.5}D^2(1+\log D)(1+\log B)A_1\cdots A_t$ is weaker,
since $1.4\cdot30^6\cdot3^{4.5}=1.4319\cdot10^{11}\ge C_1(3)$; using it would give $\ell_1=45.31$ in place of
$45.28$.

\subsection{Two smooth terms}

\begin{proposition}\label{prop:B}
In the setting of Section~\ref{sec:reduction}, assume \textup{(H)}, $A>1$ and $Y\le n$, and put $L:=\log n$.
Then
\begin{equation}\label{eq:41}
L-\log Y< C_1(3)\,(\log Y)^2\,\max(T,0.16)\,(1+\log\Bbar),\qquad\text{where }\ \Bbar:=\frac{\log(n+Y)}{\log 2}.
\end{equation}
\end{proposition}

\begin{proof}
Take $j_1\ne j_2$ from Lemma~\ref{lem:R}(v) and write $W_i:=W_{j_i}$, $q_i:=q_{j_i}$, $e_i:=e_{j_i}$. Then
$n+j_i=W_iq_i^{e_i}$, $q_1\ne q_2$ (Lemma~\ref{lem:R}(iv)), $q_i\le Y$, $e_i\ge1$ and $\log W_i\le T$ for
$i=1,2$. Define, with real logarithms,
\[
\Lambda:=\log(n+j_1)-\log(n+j_2)=e_1\log q_1-e_2\log q_2+\log(W_1/W_2).
\]
Then $\Lambda\ne0$ because $j_1\ne j_2$. By the mean value theorem $|\Lambda|=|j_1-j_2|/\xi$ for some $\xi$
between $n+j_1$ and $n+j_2$, so $\xi>n$; and $|j_1-j_2|\le Y-1-l<Y$. Hence $|\Lambda|<Y/n$, that is,
\begin{equation}\label{eq:42}
-\log|\Lambda|>L-\log Y.
\end{equation}
Since $q_i^{e_i}\le n+j_i<n+Y$, we have $e_i\le\log(n+Y)/\log q_i\le\Bbar$, so
$B^*:=\max(e_1,e_2,1)\le\Bbar$; note $\Bbar\ge1$.

\emph{Case 1: $W_1\ne W_2$.} Apply Theorem~\ref{thm:matveev} with $t=3$, $\alpha_1=q_1$, $\alpha_2=q_2$,
$\alpha_3=W_1/W_2$, $b_1=e_1$, $b_2=-e_2$, $b_3=1$. The logarithms are non-zero since $q_i\ge2$ and
$W_1/W_2\ne1$. We have $h(q_i)=\log q_i$. Writing $W_1/W_2=a/b$ in lowest terms, $a\mid W_1$ and $b\mid W_2$, so
$h(W_1/W_2)\le\log\max(W_1,W_2)\le T$; also $|\log(W_1/W_2)|\le\max(\log W_1,\log W_2)\le T$ because both
logarithms are nonnegative. So $A_1:=\log q_1$, $A_2:=\log q_2$ and $A_3:=\max(T,0.16)$ satisfy (2.4)
($\log q_i\ge\log2>0.16$). Since $\Lambda\ne0$, Theorem~\ref{thm:matveev} with $B^*$ in place of $B$ gives
\[
-\log|\Lambda|<C_1(3)\log q_1\log q_2\max(T,0.16)(1+\log B^*)\le C_1(3)(\log Y)^2\max(T,0.16)(1+\log\Bbar).
\]

\emph{Case 2: $W_1=W_2$.} Then $\ln\alpha_3=0$ is not an admissible value, and we use two logarithms instead:
$\Lambda=e_1\log q_1-e_2\log q_2\ne0$. Theorem~\ref{thm:matveev} with $t=2$, $\alpha_i=q_i$, $A_i=\log q_i$ gives
\[
-\log|\Lambda|<C_1(2)\log q_1\log q_2(1+\log B^*)\le 0.16\,C_1(3)(\log Y)^2(1+\log\Bbar),
\]
using $C_1(2)=7.47\cdot10^8\le0.16\,C_1(3)=2.22\cdot10^{10}$. This is at most the bound of Case~1.

In both cases, combining with \eqref{eq:42} gives \eqref{eq:41}.
\end{proof}

\section{The numerical chain and \texorpdfstring{$N_0$}{N0}}\label{sec:numerics}

Besides Theorem~\ref{thm:matveev} we use two explicit estimates: $x/\log x<\pi(x)$ for real $x\ge17$
(Rosser--Schoenfeld \cite[Corollary~1, (3.5)]{RS62}), and $\omega(m)\le R\log m/\log\log m$ for $m\ge3$, where
$R:=1.3841$ (Robin \cite[Th\'eor\`eme~11]{Ro83}).

\subsection{Constants}
Fix an integer $l\ge1$ and put $c:=10$. For real $\ell_1$ define
\begin{equation}\label{eq:consts}
\begin{aligned}
m&:=\log(cl), & \kappa_1&:=1+\frac{m}{\ell_1}, & \kappa_2&:=1+\frac{1-\log\log2+e^{-\ell_1}}{\ell_1},\\
\mu&:=\frac{2Rm}{c-2R}, & \rho&:=\frac{\ell_1+m-1}{\ell_1-\mu}, & \eps_1&:=\frac{3e^{-\ell_1}}{c(\ell_1+m-1)},\\
\beta&:=2c+\frac{2c(m-1)+2R/\ell_1+6e^{-\ell_1}+1}{\ell_1}, &
\eps_2&:=\frac{\beta\,\kappa_1\,\ell_1^2e^{-\ell_1}}{(c-2R)(1-\mu/\ell_1)}, &
\kappa_3&:=\frac{c}{c-2R}\,\rho\,\kappa_1\,\frac{1+\eps_1}{1-\eps_2},\\
K&:=C_1(3)\,\kappa_1^2\kappa_2\kappa_3+\frac{\kappa_1}{\ell_1^4}, &
\varphi(\ell_1)&:=\ell_1-5\log\ell_1-\log K.
\end{aligned}
\end{equation}
Here $-\log\log2=0.36651\ldots$. Note that $m\ge\log10>1$, and that $\kappa_3>1$ whenever $0<\eps_2<1$ and
$\ell_1>\mu$.

\begin{proposition}\label{prop:52}
Let $l\ge1$, and let $\ell_1\ge10$ satisfy $\ell_1>\mu$, $\eps_2<1$ and $\varphi(\ell_1)\ge0$. Then every integer $n$
with $\log\log n\ge\ell_1$ and $\log n\ge l$ has an integer $k$ with $l\le k\le 10\,l\log n$ and $v(n,k)\ge2$.
\end{proposition}

\begin{proof}
Let $L:=\log n$, $\ell:=\log L\ge\ell_1\ge10$, $y:=clL$ and $Y:=\lfloor y\rfloor$. Suppose, for contradiction,
that (H) holds: $v(n,k)\le1$ for every integer $k\in[l,y]$. Then $L\ge e^{10}$, so $y\ge10L\ge67$;
$Y>y-1\ge l+2$; $Y\le y\le cL^2<e^L=n$ (as $l\le L$); $n\ge26$ and $l-1<n$; and $\log y=\ell+m$.

\emph{S1 (primes).} By \cite[(3.5)]{RS62} with $x=y$: $\pi(y)>y/\log y=clL/(\ell+m)$.

\emph{S2 ($\omega$).} Every prime dividing $\Pi_l(n)$ divides some $n+i$ with $0\le i\le l-1$, and $n+i\ge3$, so
by \cite{Ro83}
\[
\omega_{l,y}(n)\le\sum_{i=0}^{l-1}\omega(n+i)\le\sum_{i=0}^{l-1}\frac{R\log(n+i)}{\log\log(n+i)}\le\frac{lR(L+1)}{\ell},
\]
using $\log(n+i)\le L+(l-1)/n\le L+1$ and $\log\log(n+i)\ge\log\log n=\ell>0$.

\emph{S3 (factorial).} Let $G(x):=(x+1)\log x-x+1$. For integers $M\ge1$,
$\log M!=\sum_{i=1}^{M-1}\log i+\log M\le\int_1^M\log t\,dt+\log M=G(M)$, and $G$ is increasing on $[1,\infty)$
since $G'(x)=\log x+1/x>0$. With $M=Y-l-1\in[1,y-l-1]$,
\[
\log((Y-l-1)!)\le G(y-l-1)=(y-l)\log(y-l-1)-(y-l-1)+1\le y\log y-y+l+2=lLc(\ell+m-1)+l+2.
\]

\emph{S4 ($A$).} By S1--S3 and the definition of $A$ (with $\log n=L$),
\[
A-1\ge\frac{clL}{\ell+m}-\frac{2lRL}{\ell}-\frac{2lR}{\ell}-2cl(\ell+m-1)-\frac{2(l+2)}{L}-1=lL\,a(\ell)-E(\ell),
\]
where
\[
a(\ell):=\frac{c}{\ell+m}-\frac{2R}{\ell}=\frac{(c-2R)(\ell-\mu)}{\ell(\ell+m)},\qquad
E(\ell):=\frac{2lR}{\ell}+2cl(\ell+m-1)+\frac{2(l+2)}{L}+1.
\]
Since $\ell\ge\ell_1>\mu$, $a(\ell)>0$. Using $1/l\le1$, $(l+2)/l\le3$, $1/L=e^{-\ell}\le e^{-\ell_1}$,
$1/\ell\le1/\ell_1$ and $m-1>0$,
\[
\frac{E(\ell)}{l}\le2c\ell+\Bigl(2c(m-1)+\frac{2R}{\ell_1}+6e^{-\ell_1}+1\Bigr)\le\beta\ell,\qquad
\frac1{a(\ell)}=\frac{\ell+m}{(c-2R)(1-\mu/\ell)}\le\frac{\kappa_1\ell}{(c-2R)(1-\mu/\ell_1)}.
\]
Hence $E(\ell)/(lL\,a(\ell))\le[\beta\kappa_1/((c-2R)(1-\mu/\ell_1))]\,\ell^2e^{-\ell}\le\eps_2$, because
$\ell^2e^{-\ell}$ is decreasing for $\ell\ge2$. Therefore $A-1\ge lL\,a(\ell)(1-\eps_2)>0$.

\emph{S5 ($T$).} By S3, $\log((Y-l-1)!)\le lLc(\ell+m-1)(1+\eps_1)$, since
$(l+2)/(lLc(\ell+m-1))\le3e^{-\ell_1}/(c(\ell_1+m-1))=\eps_1$. With S4,
\[
T\le\frac{c(\ell+m-1)(1+\eps_1)}{a(\ell)(1-\eps_2)}
=\frac{c}{c-2R}\cdot\frac{\ell+m-1}{\ell-\mu}\cdot\ell(\ell+m)\cdot\frac{1+\eps_1}{1-\eps_2}\le\kappa_3\ell^2,
\]
because $(\ell+m-1)/(\ell-\mu)$ is decreasing in $\ell$ (its derivative has the sign of $-(m-1+\mu)<0$), hence
at most $\rho$, and $\ell+m\le\kappa_1\ell$. Also $0.16\le\kappa_3\ell^2$, so $\max(T,0.16)\le\kappa_3\ell^2$.

\emph{S6 ($\Bbar$ and $\log Y$).} We have $\log(n+Y)\le L+Y/n\le L+1$ and $\log(L+1)\le\ell+1/L$, so
\[
1+\log\Bbar\le1+\ell+e^{-\ell_1}-\log\log2\le\kappa_2\ell,\qquad \log Y\le\log y=\ell+m\le\kappa_1\ell.
\]

\emph{S7 (conclusion).} All hypotheses of Proposition~\ref{prop:B} hold ($A>1$ by S4, and $Y\le n$). By
\eqref{eq:41} and S5--S6,
\[
L<\kappa_1\ell+C_1(3)\,\kappa_1^2\ell^2\cdot\kappa_3\ell^2\cdot\kappa_2\ell\le K\ell^5,
\]
using $\kappa_1\ell\le(\kappa_1/\ell_1^4)\ell^5$. That is, $\log n<K(\log\log n)^5$. But
$\psi(\ell):=\ell-5\log\ell$ has $\psi'(\ell)=1-5/\ell>0$ for $\ell>5$, so $\psi(\ell)\ge\psi(\ell_1)\ge\log K$,
which means $L=e^\ell\ge K\ell^5$. This contradiction proves the proposition.
\end{proof}

\subsection{Proof of Theorem A}
For $l=1$ and $\ell_1=45.28$ the script \texttt{n0\_compute.py} evaluates every quantity in \eqref{eq:consts} as
an interval (mpmath interval arithmetic, 60 digits, outward rounding) and certifies from the interval endpoints
every numerical inequality used above: $\varphi(\ell_1)\ge0$, $\ell_1\ge10$, $\ell_1>\mu$, $\eps_2<1$,
$\kappa_3\ell_1^2\ge0.16$, $m>1$, $y\ge67$ and $C_1(2)\le0.16\,C_1(3)$. The decimal inputs $1.3841$ and $45.28$
enter as enclosing intervals. The monotonicity facts are proved in the text. Table~\ref{tab:consts} lists the
values, rounded for display.

\begin{table}[ht]
\centering
\small
\begin{tabular}{@{}lll@{\qquad}lll@{}}
\toprule
quantity & value & script name & quantity & value & script name\\
\midrule
$m=\log10$ & $2.30258509299$ & \texttt{m} & $\beta$ & $20.5987817363$ & \texttt{e1}\\
$\kappa_1$ & $1.05085214428$ & \texttt{kappa1} & $\eps_2$ & $1.35404300089\cdot10^{-16}$ & \texttt{eps2}\\
$\kappa_2$ & $1.03017917227$ & \texttt{kappa2} & $\kappa_3$ & $1.52457710226$ & \texttt{kappa3}\\
$\mu$ & $0.881387214031$ & \texttt{mu} & $C_1(3)$ & $1.39007316922\cdot10^{11}$ & \texttt{C1(3)}\\
$\rho$ & $1.04919010235$ & \texttt{ratio} & $K$ & $2.41092138023\cdot10^{11}$ & \texttt{K}\\
$\eps_1$ & $1.39329644231\cdot10^{-22}$ & \texttt{eps1} & $\varphi(45.28)$ & $0.00722781770539>0$ & \texttt{phi}\\
\bottomrule
\end{tabular}
\caption{The constants \eqref{eq:consts} for $l=1$, $c=10$, $\ell_1=45.28$.}\label{tab:consts}
\end{table}

So Proposition~\ref{prop:52} applies with $l=1$; the condition $\log n\ge1$ is automatic. Every $n$ with
$\log\log n\ge45.28$, that is every $n\ge N_0=\exp(e^{45.28})$, has some $k$ with $1\le k\le10\log n$ and
$v(n,k)\ge2$. This proves Theorem~A. \qed

\medskip
The value $45.28$ is the least multiple of $0.01$ passing the certificate; $45.27$ fails. The simplifications in
S4--S7 cost less than $0.01$ in $\ell_1$: evaluated directly from S1--S3 and S6, the right-hand side of
\eqref{eq:41} first falls below $L$ near $\ell\approx45.272$ (not used), and it is at most $K\ell^5$ at
$\ell\in\{45.28,46,50,60,100,1000\}$, a check of the algebra in S4--S7.

\section{The version uniform in \texorpdfstring{$l$}{l}}\label{sec:uniform}

For fixed $l$, Proposition~\ref{prop:52} with the constants \eqref{eq:consts} gives the thresholds in
Table~\ref{tab:perl}, which proves Theorem~A$_l$. They depend on $l$ only through $m=\log(10l)$ and grow slowly:
$\ell_1(l)=45.95$, $46.53$, $47.54$ at $l=10^3,10^6,10^{12}$. For all $1\le l\le\log n$ at once,
Proposition~\ref{prop:unif} gives the single threshold $\log\log n\ge48.82$.

\begin{table}[ht]
\centering
\small
\setlength{\tabcolsep}{4.5pt}
\begin{tabular}{@{}l*{10}{c}@{}}
\toprule
$l$ & 1 & 2 & 3 & 4 & 5 & 6 & 7 & 8 & 9 & 10\\
$\ell_1(l)$ & 45.28 & 45.35 & 45.39 & 45.42 & 45.44 & 45.46 & 45.48 & 45.49 & 45.50 & 45.51\\
\midrule
$l$ & 11 & 12 & 13 & 14 & 15 & 16 & 17 & 18 & 19 & 20\\
$\ell_1(l)$ & 45.52 & 45.53 & 45.54 & 45.54 & 45.55 & 45.56 & 45.56 & 45.57 & 45.57 & 45.58\\
\bottomrule
\end{tabular}
\caption{Certified thresholds for fixed $l$ ($c=10$): $\ell_1(l)$ is the least multiple of $0.01$ passing the
certificate of Proposition~\ref{prop:52}. Correspondingly $\log N(l)=e^{\ell_1(l)}$ runs from $4.622\cdot10^{19}$
($l=1$) to $6.239\cdot10^{19}$ ($l=20$), and $K(l)$ from $2.411\cdot10^{11}$ to $3.154\cdot10^{11}$.}\label{tab:perl}
\end{table}

For the uniform version put $c=10$, $R=1.3841$, let $\kappa_2$ be as in \eqref{eq:consts}, and define
\begin{equation}\label{eq:unif}
\begin{gathered}
\kappa_1^\unif:=2+\frac{\log c}{\ell_1},\qquad \mu^\unif:=\frac{2R\log c}{c-4R},\qquad
\rho^\unif:=\frac{2\ell_1+\log c-1}{\ell_1-\mu^\unif},\qquad \eps_1^\unif:=\frac{3e^{-\ell_1}}{c(\ell_1+\log c-1)},\\
\beta^\unif:=4c+\frac{2c(\log c-1)+2R/\ell_1+6e^{-\ell_1}+1}{\ell_1},\qquad
\eps_2^\unif:=\frac{\beta^\unif\kappa_1^\unif\ell_1^2e^{-\ell_1}}{(c-4R)(1-\mu^\unif/\ell_1)},\\
\kappa_3^\unif:=\frac{c}{c-4R}\,\rho^\unif\kappa_1^\unif\,\frac{1+\eps_1^\unif}{1-\eps_2^\unif},\qquad
K^\unif:=C_1(3)(\kappa_1^\unif)^2\kappa_2\kappa_3^\unif+\kappa_1^\unif/\ell_1^4,\qquad
\varphi^\unif(\ell_1):=\ell_1-5\log\ell_1-\log K^\unif.
\end{gathered}
\end{equation}

\begin{proposition}\label{prop:unif}
Suppose $c>4R$, $\ell_1\ge10$, $\ell_1>\mu^\unif$, $\eps_2^\unif<1$ and $\varphi^\unif(\ell_1)\ge0$. Then for every
integer $l\ge1$ and every integer $n$ with $\log\log n\ge\ell_1$ and $\log n\ge l$, some $k$ with
$l\le k\le c\,l\log n$ has $v(n,k)\ge2$.
\end{proposition}

\begin{proof}
Repeat the proof of Proposition~\ref{prop:52}. Now $m=\log(cl)$ is not fixed; from $1\le l\le L$ we only know
$\log c\le m\le\ell+\log c$. The preliminary facts and S1--S3 are unchanged. The later steps change as follows.

\emph{S4.} $a(\ell)=c/(\ell+m)-2R/\ell$ is decreasing in $m$, so
$a(\ell)\ge a^\unif(\ell):=c/(2\ell+\log c)-2R/\ell=(c-4R)(\ell-\mu^\unif)/(\ell(2\ell+\log c))>0$. Using
$m\le\ell+\log c$, $E(\ell)/l\le2R/\ell_1+2c(2\ell+\log c-1)+6e^{-\ell_1}+1\le\beta^\unif\ell$ (as $\log c>1$).
Also $1/a^\unif(\ell)\le\kappa_1^\unif\ell/((c-4R)(1-\mu^\unif/\ell_1))$. Hence
$E/(lLa)\le E/(lLa^\unif)\le\eps_2^\unif$ and $A-1\ge lL\,a(\ell)(1-\eps_2^\unif)>0$.

\emph{S5.} Since $m\ge\log c$, $(l+2)/(lLc(\ell+m-1))\le\eps_1^\unif$. Also
$c(\ell+m-1)/a(\ell)\le c(2\ell+\log c-1)/a^\unif(\ell)\le\frac{c}{c-4R}\rho^\unif\kappa_1^\unif\ell^2$, because
$(2\ell+\log c-1)/(\ell-\mu^\unif)$ is decreasing (its derivative has the sign of $-(2\mu^\unif+\log c-1)<0$).
So $T\le\kappa_3^\unif\ell^2$, and $\kappa_3^\unif\ell_1^2\ge0.16$.

\emph{S6.} $\log Y\le\ell+m\le2\ell+\log c\le\kappa_1^\unif\ell$. The bound $1+\log\Bbar\le\kappa_2\ell$ is
unchanged; it uses $Y\le cL^2<n$.

\emph{S7.} $L<K^\unif\ell^5$, which contradicts $\varphi^\unif(\ell_1)\ge0$ as before.
\end{proof}

With $c=10$ (and $4R=5.5364<10$) the script certifies $\ell_1=48.82$ (and $48.81$ fails), with
$\kappa_1^\unif=2.04716$, $\mu^\unif=1.42800$, $\rho^\unif=2.08775$, $\beta^\unif=40.5553$,
$\eps_2^\unif=2.87\cdot10^{-17}$, $\kappa_3^\unif=9.57515$, $K^\unif=5.73427\cdot10^{12}$ and
$\varphi^\unif(48.82)=0.00182>0$; it also certifies $\ell_1>\mu^\unif$, $\eps_2^\unif<1$,
$\kappa_3^\unif\ell_1^2\ge0.16$ and $\log c>1$. This proves Theorem~A$_{\mathrm{unif}}$: if $n\ge N(l)$ then $\log\log n\ge48.82$ and $\log n\ge l$.
Corollary~A$_l$ follows as Corollary~A1 does. \qed

\section{Remarks}\label{sec:remarks}

\subsection{The Langevin route}\label{sec:langevin}
The qualitative content of Corollary~A1 follows from a theorem of Langevin and the case $l=1$ of
Lemma~\ref{lem:R}(i), restated here; this route was pointed out in the prior-art check for this note.

\begin{lemma}\label{lem:G2}
Let $n\ge1$ and $K\ge2$ be integers, and assume $v(n,k)\le1$ for every $1\le k\le K$. For each prime $p\le K$ with
$p\nmid n$ let $k_0(p)$ be the least $k\ge1$ with $p\mid n+k$. Then $1\le k_0(p)\le p-1$, $P(n+k_0(p))=p$, and
$p\mapsto k_0(p)$ is injective. Hence $\#\{1\le k\le K: P(n+k)\le K\}\ge\pi(K)-\omega(n)$.
\end{lemma}

\begin{proof}
Since $p\nmid n$, one of $n+1,\dots,n+p-1$ is divisible by $p$, so $1\le k_0(p)\le p-1<p$ and $p$ is a prime
factor of $n+k_0(p)$ greater than $k_0(p)$. As $v(n,k_0(p))\le1$, Lemma~\ref{lem:v} shows that $p$ is the only
such factor; every other prime factor of $n+k_0(p)$ is at most $k_0(p)<p$, so $P(n+k_0(p))=p\le K$. This
recovers $p$ from $k_0(p)$, so the map is injective. The number of primes $p\le K$ with $p\nmid n$ is
$\pi(K)-\#\{p\le K:p\mid n\}$, which is at least $\pi(K)-\omega(n)$. By injectivity these yield at least
$\pi(K)-\omega(n)$ distinct $k\in[1,K]$ with $P(n+k)\le K$.
\end{proof}

Langevin \cite[Corollaire~1]{La81} considers integers $n,k,a$ with $n>1$, $a\ge1$, $(n,a)=1$, $k\ge2$ and
$a<n^{t_1}$ (his condition (1)). His constants $C$ and $c$ are effectively computable in terms of the parameters
$t_1,\dots,t_6$ of his Th\'eor\`eme~1, which allows any real $t_1<1$; his Exemple (p.~242) takes $t_1=\frac12$,
$t_2=1$, $t_3=\frac13$, $t_4=3$, $t_5=t_6=6$. Part (8) of Corollaire~1 states: for every $\eps>0$ there is an
explicit $r(\eps)$, depending on $C$, $c$ and $\eps$, such that if $\log n/\log k\ge r(\eps)$ then
\[
\operatorname{card}\cY<\max\Bigl(2,\ (1+c+\eps)\,\pi(k)\,\frac{\log\log X}{\log X}\Bigr),\qquad
X:=\frac{\log n}{(\log k)^2},
\]
where $\cY$ is the set of the numbers $n+ia$ $(1\le i\le k)$ with $P(n+ia)\le k$.

\begin{proposition}\label{prop:langevin}
Assuming \cite[Corollaire~1, (8)]{La81}, there is an effectively computable $N_1$ such that $v_1(n)\ge2$ for all
$n\ge N_1$; in particular $\{n:v_1(n)=1\}$ is finite.
\end{proposition}

\begin{proof}
Fix Langevin's example parameters (so $t_1=\frac12$) and $\eps=1$. For large $n$ put $k^*:=\lfloor10\log n\rfloor$
and $a=1$; condition (1) holds since $1<n^{1/2}$. As $n\to\infty$, $\log n/\log k^*\to\infty$ and
$X=\log n/(\log k^*)^2\to\infty$, so for large $n$ (8) applies and gives
$\operatorname{card}\cY<\max(2,o(\pi(k^*)))$. Suppose $v(n,k)\le1$ for every $1\le k\le k^*$. Lemma~\ref{lem:G2}
with $K=k^*$ gives $\operatorname{card}\cY\ge\pi(k^*)-\omega(n)$. By \cite{Ro83}, $\omega(n)\le R\log n/\log\log n$,
and by \cite{RS62}, $\pi(k^*)>k^*/\log k^*$ with $k^*>10\log n-1$ and $\log k^*\le\log\log n+\log10$; hence
$\omega(n)\le(R/10+o(1))\,\pi(k^*)$ and $\operatorname{card}\cY\ge(1-0.13841-o(1))\,\pi(k^*)$, which contradicts the
upper bound for large $n$. So for all $n\ge N_1$ some $k\le10\log n$ has $v(n,k)\ge2$, and every constant involved
is effective.
\end{proof}

With Lemma~\ref{lem:R}(i) in place of Lemma~\ref{lem:G2} and $k^*:=\lfloor10\,l\log n\rfloor$, the same argument
gives Corollary~A$_l$ qualitatively, also uniformly for $1\le l\le\log n$: the numbers $n+j$ $(j\in J)$ lie in
$\cY$, $\omega_{l,k^*}(n)\le(2R/10+o(1))\pi(k^*)$ and $X\ge L/(2\ell+\log10)^2\to\infty$.

We have not refereed Langevin's proof; we rely on the statement as printed. Langevin's constants are effective. We have not computed a certified numerical onset for this alternative
route and use it only to establish the qualitative consequence.
By contrast, the bound of Ramachandra, Shorey and Tijdeman \cite[Theorem~2]{RST76}, that at most $\pi(k)$ of
$u+1,\dots,u+k$ are $k$-smooth when $u\ge\exp(C(\log k)^2)$, is exactly borderline against Lemma~\ref{lem:G2} and
gives no contradiction by itself.

\subsection{Sharper inputs and the constant 10}
Three sharper inputs lower the threshold of Theorem~A to $\log\log n\ge44.96$, that is
$\log N_0'=e^{44.96}\approx3.356\cdot10^{19}$, a factor $0.73$ below $\log N_0$: in S1 the bound
$\pi(y)>y/(\log y-\frac12)$ for $y\ge67$ \cite[(3.3)]{RS62}; in S2 Robin's
$\omega(m')\le\log m'/(\log\log m'-1.1714)$ for $m'\ge26$ \cite[Th\'eor\`eme~13]{Ro83}; and in Case~1 of
Proposition~\ref{prop:B} Matveev's $B$ from (1.3) in place of $B^*$, with $\alpha_3=W_1/W_2$ ordered last and
$A_3:=\kappa_3'\ell^2\ge T$. This gives $L<\kappa_1\ell+C_1(3)\kappa_1^2\kappa_3'\ell^4(\ell-d')$ with explicit
constants ($\kappa_3'=1.39730$ and $d'=6.94609$ at $\ell_1=44.96$). The function
$Q(\ell):=(\kappa_1\ell+C_1(3)\kappa_1^2\kappa_3'\ell^4(\ell-d'))e^{-\ell}$ is non-increasing for $\ell\ge\ell_1$
(its log-derivative is at most $4/\ell+1/(\ell-d')-1<0$), and the script certifies $Q(44.96)=0.99329<1$ with all
side conditions, including that the two-logarithm case is dominated; $44.95$ fails. The derivation is written out
in the accompanying notes (\texttt{PROOF\_THEOREM\_A.md}, Section~7.2). We state Theorem~A with the simpler $45.28$.

The constant $10$ is not essential. For fixed $l$ the argument works with any $c>2R=2.7682$, since then
$a(\ell)>0$ for large $\ell$, with a threshold that grows as $c$ decreases towards $2R$; with Robin's
Th\'eor\`eme~13 any $c>2$ works. The uniform version as written needs $c>4R$. We have not tried to optimise $c$ or
the thresholds.

\subsection{Finite data}\label{sec:finite}
These computations are evidence only; none is used in the proofs.
\emph{Lemma~\ref{lem:R}.} The script's check mode verifies every conclusion of Lemma~\ref{lem:R}(i)--(v) on every
triple $(n,l,Y)$ with $3\le n\le10^6$, $l\in\{1,2,3\}$, $l+2\le Y\le150$ and $v(n,k)\le1$ for all $l\le k\le Y$:
$91{,}709$ triples, $40{,}312$ of them with $|J_1|\ge2$, and no failure. The referee's separate implementation in
exact integer arithmetic (not distributed) checked $207{,}874$ instances with $l\le4$ and $Y\le200$, also
without failure.

\emph{Least witnesses up to $10^7$.} For every $1\le n\le10^7$ the least $k\ge1$ with $v(n,k)\ge2$ was searched
in the window $k\le\lfloor30\log(n+2)+30\rfloor$, which contains $[1,10\log n]$. Exactly $29$ values of $n$ have no
such $k$ in the window, namely the $29$ values of Erd\H{o}s and Selfridge listed in the introduction; so for
$2\le n\le10^7$ the conclusion of Theorem~A fails exactly at the $28$ values $2\le n\le330$ of that list. For every
other $n\le10^7$ the least such $k$ is at most $11$, and $k=11$ occurs only at $n=210$ ($221=13\cdot17$) and
$n=840$ ($851=23\cdot37$).

\emph{Exception lists for $l\le20$.} For $l=1,2,3,5,8,13,20$ the numbers $n\le10^7$ with no
$k\in[l,\lfloor40\log(n+2)+40\rfloor]$ satisfying $v(n,k)\ge2$ number $29$, $103$, $187$, $383$, $722$, $1531$,
$2817$, with largest elements $330$, $1365$, $1365$, $2415$, $4895$, $11655$, $27714$. This truncates
$\sup_{k\ge l}$: every $n$ not listed has $v_l(n)\ge2$, but a listed $n$ is only known to have no witness in the
window.

The witness computation was cross-checked against an independent brute-force program for $n\le2\cdot10^6$,
and the two rows with least $k>10$ were re-verified, including minimality of $k$, by an independently written
trial-division checker sharing no code with the generator.

\subsection{Open questions not addressed}\label{sec:open}
This note does not address the main problem $v_0(n)\to\infty$, nor $v_l(n)\to\infty$ for fixed $l$; nor
$v_0(n)\ge3$ or $v_l(n)\ge3$ for all large $n$; nor the conjectures that $330$ is the largest $n$ with $v_1(n)=1$
and that $80$ is the largest $n$ with $V_1(n)=1$. Since $N_0$ is far beyond computation, the set
$\{n:v_1(n)=1\}$ is still not determined.

\section*{Provenance and verification status}
The proof in this note was found and written by an AI pipeline built on Claude models (Anthropic): Claude Opus
for proof search and writing, an independent Claude referee working in a fresh context, and Claude Fable
coordinating. The referee checked every step against the primary sources and recomputed the constants
independently; its verdict was ``pass with repairs'', the repairs concerned wording and citations only, and they
are incorporated here. A subsequent cross-vendor review by OpenAI Codex on 26 September 2026 gave
PASS-WITH-REPAIRS with $N_0$ unchanged; its repairs are incorporated here (see
\texttt{REFEREE\_ASTRA\_20260926.md}). No human referee or kernel formalization is claimed.
The numerical constants were certified by interval arithmetic (\texttt{n0\_compute.py}, mpmath); the witness
computation was checked by an independently written program. The author takes responsibility for the
mathematics. There is no Lean formalisation; the proof uses Matveev's theorem, which is not available in Mathlib.
The literature search was done by the same pipeline, in two passes without general web search. The following
sources could not be consulted: Erd\H{o}s \cite[p.~178]{Er98}; Shorey and Tijdeman \cite{ST86}; Erd\H{o}s
\cite{Er76}; Tijdeman's review in Math.\ Rev.\ 54 (1977), cited in \cite{La81} for the case $a=1$ of (8); and the
Math.\ Rev.\ review of \cite{La81}. Any of them may already contain Theorem~A or its qualitative form. Code and
data: \url{https://github.com/chy4pro/automath}, directory \texttt{problems/erdos889}.

\begin{thebibliography}{RST76}\footnotesize\raggedright\setlength{\itemsep}{0.5pt}
\bibitem[Er76]{Er76} P.~Erd\H{o}s, Problems and results on consecutive integers, \emph{Publ.\ Math.\ Debrecen}
\textbf{23} (1976), 271--282.
\bibitem[Er98]{Er98} P.~Erd\H{o}s, Some of my new and almost new problems and results in combinatorial number
theory, in: \emph{Number Theory (Eger, 1996)}, de Gruyter, Berlin, 1998, 169--180.
\bibitem[ErSe67]{ErSe67} P.~Erd\H{o}s and J.~L.~Selfridge, Some problems on the prime factors of consecutive
integers, \emph{Illinois J.\ Math.} \textbf{11} (1967), 428--430.
\bibitem[EP889]{EP889} T.~F.~Bloom, Erd\H{o}s Problem \#889, \url{https://www.erdosproblems.com/889},
accessed 2026-09-25.
\bibitem[FC]{FC} Google DeepMind, Formal Conjectures, file \texttt{FormalConjectures/ErdosProblems/889.lean},
\url{https://github.com/google-deepmind/formal-conjectures}, accessed 2026-09-25.
\bibitem[Gu04]{Gu04} R.~K.~Guy, \emph{Unsolved Problems in Number Theory}, 3rd ed., Problem Books in
Mathematics, Springer, New York, 2004, Problem B27.
\bibitem[Ki26]{Ki26} K.~Kitamura, \emph{A Lean 4 proof that the capital-$V_1$ exceptional set in Erd\H{o}s
Problem 889 is finite}, GitHub repository \texttt{KitaKen1/erdos-889-capital-v1-finite}, created 2026-09-22,
\url{https://github.com/KitaKen1/erdos-889-capital-v1-finite}, accessed 2026-09-25.
\bibitem[La81]{La81} M.~Langevin, Facteurs premiers d'entiers en progression arithm\'etique, \emph{Acta Arith.}
\textbf{39} (1981), 241--249.
\bibitem[Ma00]{Ma00} E.~M.~Matveev, An explicit lower bound for a homogeneous rational linear form in the
logarithms of algebraic numbers.~II, \emph{Izv.\ Ross.\ Akad.\ Nauk Ser.\ Mat.} \textbf{64} (2000), no.~6,
125--180; English translation in \emph{Izv.\ Math.} \textbf{64} (2000), no.~6, 1217--1269.
\bibitem[RST75]{RST75} K.~Ramachandra, T.~N.~Shorey and R.~Tijdeman, On Grimm's problem relating to
factorisation of a block of consecutive integers, \emph{J.\ reine angew.\ Math.} \textbf{273} (1975), 109--124.
\bibitem[RST76]{RST76} K.~Ramachandra, T.~N.~Shorey and R.~Tijdeman, On Grimm's problem relating to
factorisation of a block of consecutive integers.~II, \emph{J.\ reine angew.\ Math.} \textbf{288} (1976),
192--201.
\bibitem[Ro83]{Ro83} G.~Robin, Estimation de la fonction de Tchebychef $\theta$ sur le $k$-i\`eme nombre premier
et grandes valeurs de la fonction $\omega(n)$ nombre de diviseurs premiers de $n$, \emph{Acta Arith.} \textbf{42}
(1983), 367--389.
\bibitem[RS62]{RS62} J.~B.~Rosser and L.~Schoenfeld, Approximate formulas for some functions of prime numbers,
\emph{Illinois J.\ Math.} \textbf{6} (1962), 64--94.
\bibitem[ST86]{ST86} T.~N.~Shorey and R.~Tijdeman, \emph{Exponential Diophantine Equations}, Cambridge Tracts in
Mathematics 87, Cambridge University Press, Cambridge, 1986.
\bibitem[Ti73]{Ti73} R.~Tijdeman, On integers with many small prime factors, \emph{Compositio Math.} \textbf{26}
(1973), 319--330.
\end{thebibliography}

\end{document}
